\begin{lemma}
\label{lemma-recover-module-from-glueing-data}
Assume $\varphi : R \to S$ is a flat ring map and
$I = (f_1, \ldots, f_t) \subset R$ is an ideal such that
$R/I \to S/IS$ is an isomorphism.
Let $M$ be an $R$-module. Then the
complex (\ref{equation-glueing-complex})
is exact.
\end{lemma}

\begin{proof}
First proof. Denote $\check{\mathcal{C}}_R \to \check{\mathcal{C}}_S$
the quasi-isomorphism of extended alternating {\v C}ech complexes of
Lemma \ref{lemma-map-identifies-koszul-and-cech-complexes}.
Since these complexes are bounded with flat terms, we see that
$M \otimes_R \check{\mathcal{C}}_R \to M \otimes_R \check{\mathcal{C}}_S$
is a quasi-isomorphism too (Lemmas
\ref{lemma-derived-tor-quasi-isomorphism} and
\ref{lemma-derived-tor-quasi-isomorphism-other-side}). Now the complex
(\ref{equation-glueing-complex}) is a truncation of the cone
of the map
$M \otimes_R \check{\mathcal{C}}_R \to M \otimes_R \check{\mathcal{C}}_S$
and we win.

\medskip\noindent
Second computational proof.
Let $m \in M$. If $\alpha(m) = 0$, then $m \in M[I^\infty]$, see
Lemma \ref{lemma-torsion-free}. Pick $n$ such that $I^n m = 0$
and consider the map $\varphi : R/I^n \to M$.
If $m \otimes 1 = 0$, then $\varphi \otimes 1_S = 0$, hence
$\varphi = 0$ (see
Lemma \ref{lemma-neighbourhood-equivalence})
hence $m = 0$. In this way we see that $\alpha$ is injective.

\medskip\noindent
Let $(m', m'_1, \ldots, m'_t) \in \Ker(\beta)$.
Write $m'_i = m_i/f_i^n$ for some $n > 0$ and $m_i \in M$.
We may, after possibly enlarging $n$ assume that
$f_i^n m' = m_i \otimes 1$ in $M \otimes_R S$ and
$f_j^nm_i - f_i^nm_j = 0$ in $M$.
In particular we see that
$(m_1, \ldots, m_t)$ defines an element $\xi$ of
$H_1(M, (f_1^n, \ldots, f_t^n))$.
Since $H_1(M, (f_1^n, \ldots, f_t^n))$ is annihilated by $I^{tn + 1}$ (see
Lemma \ref{lemma-koszul-homology-annihilated})
and since $R \to S$ is flat we see that
$$
H_1(M, (f_1^n, \ldots, f_t^n)) =
H_1(M, (f_1^n, \ldots, f_t^n)) \otimes_R S =
H_1(M \otimes_R S, (f_1^n, \ldots, f_t^n))
$$
by
Lemma \ref{lemma-neighbourhood-isomorphism}
The existence of $m'$ implies that $\xi$ maps to zero in the last group, i.e.,
the element $\xi$ is zero. Thus there exists an $m \in M$ such that
$m_i = f_i^n m$. Then $(m', m'_1, \ldots, m'_t) - \alpha(m)
= (m'', 0, \ldots, 0)$ for some $m'' \in (M \otimes_R S)[(IS)^\infty]$.
By
Lemma \ref{lemma-neighbourhood-equivalence}
we conclude that $m'' \in M[I^\infty]$ and we win.
\end{proof}